% GENERATED FILE. Edit canonical lessons, then run npm run generate:books.
% canonical-source-hash: fnv1a64:915126f10dcd3e22
% canonical-lessons: LA-C389-servitus, LA-C389-manumitto, LA-C389-pietas, LA-C389-onus, LA-C389-obsequium

\chapter{Duty and Freedom}
\label{ch:la-duty-and-freedom}
\hlchaptermodality{389}

\begin{chapteropening}
\textbf{By the end of this chapter:} \emph{I can talk about slavery and setting a slave free, about duty to the gods and to family, a burden and obedience.}

Complete the last lesson of chapter 389: I can talk about slavery and setting a slave free, about duty to the gods and to family, a burden and obedience.
\end{chapteropening}

\section[servitūs, servitūtis]{servitūs, servitūtis — slavery, servitude}

\label{lesson:LA-C389-servitus}



\begin{quote}
Before the new one: say the Latin for a companion, then the Latin for an association.
\end{quote}

\subsection*{You'll want to know: servitūs, servitūtis}
\textbf{servitūs, servitūtis} — \textquotedblleft{}slavery, servitude\textquotedblright{}.

\begin{quote}
Lībertās servitūte melior est.
\end{quote}

\begin{quote}
\emph{Freedom is better than slavery.}
\end{quote}

\emph{servitūs, servitūtis} f., built on \emph{servus}, a slave. English \emph{servitude}.

\begin{cousinweb}
\emph{servitūs} is built on \emph{servus} \textquotedblleft{}a slave\textquotedblright{} (EtymDB; Lewis \& Short), with the ending \emph{-tūs} that makes a state.
\end{cousinweb}

\begin{grammarlens}[title={What you've built}]
One more word for telling, explaining and describing at B1.
\end{grammarlens}

\subsection*{Your turn}
\begin{itemize}
\raggedright
  \item \emph{Say it:} \emph{servitūs, servitūtis}
  \item \emph{Say it:} \emph{servitūs, servitūtis}, once more
  \item \emph{Say it:} a sentence of your own with \emph{servitūs, servitūtis}
  \item \emph{Recall:} say the Latin for a companion, then the Latin for an association, then say \emph{servitūs, servitūtis} again
\end{itemize}

\begin{tcolorbox}[breakable,colback=teal!4,colframe=teal!35!black,title={Before you move on}]
What does servitūs, servitūtis mean? (\textquotedblleft{}Slavery, servitude\textquotedblright{}.) Say it once more.
\end{tcolorbox}
\section[manūmittō, manūmittere]{manūmittō, manūmittere — to set free (a slave)}

\label{lesson:LA-C389-manumitto}



\begin{quote}
Before the new one: say the Latin for an association, then the Latin for slavery.
\end{quote}

\subsection*{You'll want to know: manūmittō, manūmittere}
\textbf{manūmittō, manūmittere} — \textquotedblleft{}to set free (a slave)\textquotedblright{}.

\begin{quote}
Dominus servum fidēlem manūmīsit.
\end{quote}

\begin{quote}
\emph{The master freed his faithful slave.}
\end{quote}

\emph{manūmittō, manūmittere, manūmīsī, manūmissum}: \emph{manū} from the hand + \emph{mittō}, to let go. The freed slave became a \emph{lībertus}. English \emph{manumit}.

\begin{cousinweb}
\emph{manūmittō} is built from \emph{manū} \textquotedblleft{}from the hand\textquotedblright{}, the ablative of \emph{manus}, and \emph{mittō} \textquotedblleft{}to let go\textquotedblright{} (Lewis \& Short): to release from one's power. English \emph{manumit} comes from it.
\end{cousinweb}

\begin{grammarlens}[title={What you've built}]
One more word for telling, explaining and describing at B1.
\end{grammarlens}

\subsection*{Your turn}
\begin{itemize}
\raggedright
  \item \emph{Say it:} \emph{manūmittō, manūmittere}
  \item \emph{Say it:} \emph{manūmittō, manūmittere}, once more
  \item \emph{Say it:} a sentence of your own with \emph{manūmittō, manūmittere}
  \item \emph{Recall:} say the Latin for an association, then the Latin for slavery, then say \emph{manūmittō, manūmittere} again
\end{itemize}

\begin{tcolorbox}[breakable,colback=teal!4,colframe=teal!35!black,title={Before you move on}]
What does manūmittō, manūmittere mean? (\textquotedblleft{}To set free (a slave)\textquotedblright{}.) Say it once more.
\end{tcolorbox}
\section[pietās, pietātis]{pietās, pietātis — dutifulness, devotion (to gods, family, country)}

\label{lesson:LA-C389-pietas}



\begin{quote}
Before the new one: say the Latin for slavery, then the Latin for to set free.
\end{quote}

\subsection*{You'll want to know: pietās, pietātis}
\textbf{pietās, pietātis} — \textquotedblleft{}dutifulness, devotion (to gods, family, country)\textquotedblright{}.

\begin{quote}
Fīlius magnā pietāte patrem cūrābat.
\end{quote}

\begin{quote}
\emph{The son looked after his father with great devotion.}
\end{quote}

\emph{pietās, pietātis} f., built on \emph{pius}, dutiful: the loyalty owed to gods, parents and country. English \emph{piety} and \emph{pity}.

\begin{cousinweb}
\emph{pietās} is built on \emph{pius} \textquotedblleft{}dutiful\textquotedblright{} (EtymDB; Lewis \& Short), with the ending \emph{-tās} that makes a quality. English \emph{piety} and \emph{pity} both come from it.
\end{cousinweb}

\begin{grammarlens}[title={What you've built}]
One more word for telling, explaining and describing at B1.
\end{grammarlens}

\subsection*{Your turn}
\begin{itemize}
\raggedright
  \item \emph{Say it:} \emph{pietās, pietātis}
  \item \emph{Say it:} \emph{pietās, pietātis}, once more
  \item \emph{Say it:} a sentence of your own with \emph{pietās, pietātis}
  \item \emph{Recall:} say the Latin for slavery, then the Latin for to set free, then say \emph{pietās, pietātis} again
\end{itemize}

\begin{tcolorbox}[breakable,colback=teal!4,colframe=teal!35!black,title={Before you move on}]
What does pietās, pietātis mean? (\textquotedblleft{}Dutifulness, devotion (to gods, family, country)\textquotedblright{}.) Say it once more.
\end{tcolorbox}
\section[onus, oneris]{onus, oneris — a load, a burden}

\label{lesson:LA-C389-onus}



\begin{quote}
Before the new one: say the Latin for to set free, then the Latin for dutifulness.
\end{quote}

\subsection*{You'll want to know: onus, oneris}
\textbf{onus, oneris} — \textquotedblleft{}a load, a burden\textquotedblright{}.

\begin{quote}
Asinus onus grave portat.
\end{quote}

\begin{quote}
\emph{The donkey carries a heavy load.}
\end{quote}

\emph{onus, oneris} n. A load on a back, and a burden of duty or tax. English \emph{onus} and \emph{onerous}.

\begin{cousinweb}
\emph{onus} is inherited from a Proto-Indo-European noun meaning \textquotedblleft{}a burden\textquotedblright{} (EtymDB). English took it over as \emph{onus}, and \emph{onerous} comes from its adjective.
\end{cousinweb}

\begin{grammarlens}[title={What you've built}]
One more word for telling, explaining and describing at B1.
\end{grammarlens}

\subsection*{Your turn}
\begin{itemize}
\raggedright
  \item \emph{Say it:} \emph{onus, oneris}
  \item \emph{Say it:} \emph{onus, oneris}, once more
  \item \emph{Say it:} a sentence of your own with \emph{onus, oneris}
  \item \emph{Recall:} say the Latin for to set free, then the Latin for dutifulness, then say \emph{onus, oneris} again
\end{itemize}

\begin{tcolorbox}[breakable,colback=teal!4,colframe=teal!35!black,title={Before you move on}]
What does onus, oneris mean? (\textquotedblleft{}A load, a burden\textquotedblright{}.) Say it once more.
\end{tcolorbox}
\section[obsequium, obsequiī]{obsequium, obsequiī — compliance, obedience}

\label{lesson:LA-C389-obsequium}



\begin{quote}
Before the new one: say the Latin for dutifulness, then the Latin for a load.
\end{quote}

\subsection*{You'll want to know: obsequium, obsequiī}
\textbf{obsequium, obsequiī} — \textquotedblleft{}compliance, obedience\textquotedblright{}.

\begin{quote}
Obsequium amīcōs, vēritās odium parit.
\end{quote}

\begin{quote}
\emph{Giving way wins friends; truth breeds hatred.}
\end{quote}

\emph{obsequium, obsequiī} n., built on \emph{obsequor}, to give way to someone. English \emph{obsequious}.

\begin{cousinweb}
\emph{obsequium} is built on \emph{obsequor} \textquotedblleft{}to comply\textquotedblright{} (EtymDB; Lewis \& Short), from \emph{ob} and \emph{sequor} \textquotedblleft{}to follow\textquotedblright{}. English \emph{obsequious} comes from its adjective.
\end{cousinweb}

\begin{grammarlens}[title={What you've built}]
One more word for telling, explaining and describing at B1.
\end{grammarlens}

\subsection*{Your turn}
\begin{itemize}
\raggedright
  \item \emph{Say it:} \emph{obsequium, obsequiī}
  \item \emph{Say it:} \emph{obsequium, obsequiī}, once more
  \item \emph{Say it:} a sentence of your own with \emph{obsequium, obsequiī}
  \item \emph{Recall:} say the Latin for dutifulness, then the Latin for a load, then say \emph{obsequium, obsequiī} again
\end{itemize}

\begin{tcolorbox}[breakable,colback=teal!4,colframe=teal!35!black,title={Before you move on}]
What does obsequium, obsequiī mean? (\textquotedblleft{}Compliance, obedience\textquotedblright{}.) Say it once more.
\end{tcolorbox}
